\section{Positive trace paths and a uniform gap bound}
\label{sec:gap}

The first task is to bound the distance, along a Markov ray, between
the single curve and the factors of the comparison word. Positivity
controls the contribution of all matrix paths except one. An algebraic
norm then prevents the coefficient of that path from agreeing exactly
with the leading coefficient of the single curve. The resulting bound
depends only on the number of occurrences of the first generator.

\begin{theorem}[The gap bound]\label{thm:gap-bound}
Let $(A,B)$ be a marked modular-torus basis with positive integral
normalized Markov traces. Let $h\geqslant2$, $k\geqslant1$, and let
$\epsilon_1,\ldots,\epsilon_h\in\{0,1\}$ have sum $t$, where
$1\leqslant t<h$. If
\begin{equation}\label{eq:gap-original-equality}
 \tr A=\tr\prod_{i=1}^h AB^{k+\epsilon_i},
\end{equation}
then
\begin{equation}\label{eq:gap-final-bound}
 k\leqslant b(h):=3h^2+3h-2+\left\lfloor\frac{3h}{2}\right\rfloor.
\end{equation}
No balance or coprimality hypothesis on the gap pattern is required
for this trace inequality. For the primitive simple-curve application,
the pattern is cyclically mechanical and $\gcd(h,t)=1$.
\end{theorem}

\subsection{The positive model and the two index orientations}

We use the minimum-ray normalization from Section~\ref{sec:rays}.
In particular,
\begin{equation}\label{eq:gap-ray-data}
 \begin{gathered}
 u_j=c\lambda^j+d\lambda^{-j},\qquad
 D=9p^2-4,\qquad \lambda=\frac{3p+\sqrt D}{2},\qquad
 K=cd=\frac{p^2}{D},\\
 \lambda^{-1}\leqslant c/d\leqslant\lambda,\qquad
 c,d<\sqrt p,\qquad
 \frac52p<\lambda<3p,\qquad
 5p^2\leqslant D<9p^2,\qquad \frac19<K\leqslant\frac15.
 \end{gathered}
\end{equation}
Here $p$ and every $u_j$ are positive integers, and $u_0$ is a minimum
over all integral indices. Real conjugation gives the trace model
\begin{equation}\label{eq:gap-positive-model}
 A_0=\begin{pmatrix}3c&2/\sqrt D\\2/\sqrt D&3d\end{pmatrix},
 \qquad B=\begin{pmatrix}\lambda&0\\0&\lambda^{-1}\end{pmatrix}.
\end{equation}
Indeed the diagonal entries of the shifted first generator are $3c,3d$.
Its off-diagonal product is $9cd-1=4/D>0$, so a diagonal conjugation
makes both off-diagonal entries equal to $2/\sqrt D$. This is a change
of matrix coordinates, not an isometry assertion. Every $A_0B^j$ is
strictly positive, even for negative $j$, and has determinant one.
Every product used below consequently has trace greater than two.

After translating the minimum to zero, let the original single generator
have index $n_0$ and the product have base index $j_0$. Then
$j_0-n_0=k$. If $j_0\geqslant0$, put $l=j_0$. If $j_0<0$, reverse
the whole ray and use
\begin{equation}\label{eq:gap-complement}
 -j_0-\epsilon_i=(-j_0-1)+(1-\epsilon_i).
\end{equation}
Thus $l=-j_0-1\geqslant0$, the gap bits are complemented, and their
sum changes from $t$ to $h-t$. Both sums lie strictly between zero and
$h$. We henceforth use $t$ for the sum in the resulting orientation.
Write the single index as $\eps m$, with $m\geqslant0$. The original
quotient is recovered by one of the two formulas
\begin{equation}\label{eq:gap-reconstruction}
 k=l-\eps m
 \quad\hbox{or}\quad k=\eps m-l-1,
 \qquad k\leqslant l+m.
\end{equation}
The second formula accompanies the complementary pattern. These
formulas retain both directions of the ray and every possible position
of its minimum.

The single index cannot be zero in \eqref{eq:gap-original-equality}.
Before reversal, that case would have $n_0=0$ and $j_0=k\geqslant1$.
The constant upper-state contribution to the product trace divided by
three is
$$
 \frac{(3c\lambda^{j_0})^h\lambda^t}{3}.
$$
The minimum ratio gives $c^2\geqslant K/\lambda$, and $9K>1$, so
$3c\lambda>\sqrt\lambda$. Since $h\geqslant2$ and $t\geqslant1$,
this contribution alone exceeds
$$
 \lambda^2/3>\frac{25}{12}p^2>2\sqrt p>c+d=u_0.
$$
Every other contribution is positive. Equality is therefore impossible,
and from now on $m\geqslant1$.

\subsection{The exact path expansion}

Put
\begin{equation}\label{eq:gap-word-trace}
 V_l=\frac13\tr\prod_{i=1}^h A_0B^{l+\epsilon_i}.
\end{equation}
A term in the trace is specified by a cyclic choice of upper and lower
states at the $h$ factor boundaries. Let $j$ be its number of lower
positions, let $2r$ be its number of transitions, and let $e$ count the
long-gap positions among its lower positions. The number of transitions
is even because the path closes. Its exact contribution is
\begin{equation}\label{eq:gap-path-term}
 \frac{(3c)^{h-j-r}(3d)^{j-r}(4/D)^r}{3}
 \lambda^{(h-2j)l+t-2e}.
\end{equation}
The exponents $h-j-r$ and $j-r$ count diagonal steps and are
nonnegative. All terms in \eqref{eq:gap-path-term} are positive.
Using $cd=p^2/D$, the same coefficient can be written as
\begin{equation}\label{eq:gap-transition-factor}
 \frac{(3c)^{h-j}(3d)^j}{3}\sigma^r,
 \qquad \sigma=\frac4{9p^2}<1.
\end{equation}
The two paths without transitions contribute
\begin{equation}\label{eq:gap-constant-paths}
 X_l=3^{h-1}c^h\lambda^{hl+t},\qquad
 X_l^-=3^{h-1}d^h\lambda^{-hl-t}.
\end{equation}
Write $V_l=X_l+\mathcal R_l$, where $\mathcal R_l>0$.

The relative contribution of a path is
\begin{equation}\label{eq:gap-relative-path}
 \sigma^r(d/c)^j\lambda^{-2jl-2e}.
\end{equation}
It is at most $\lambda^{j(1-2l)}$. There are $\binom hj$ indexed
state paths with $j$ lower positions, irrespective of their transition
counts. Summing gives
\begin{equation}\label{eq:gap-binomial-bound}
 1<\frac{V_l}{X_l}\leqslant
 \begin{cases}
 (1+\lambda^{-1})^h,&l\geqslant1,\\
 (1+\lambda)^h,&l=0.
 \end{cases}
\end{equation}
This use of the binomial theorem is an analytic sum for arbitrary $h$;
it does not require listing the $2^h$ paths.

We shall also need an absolute remainder estimate. Every path in
$\mathcal R_l$ has at least one lower position, so its exponent in
\eqref{eq:gap-path-term} is at most $(h-2)l+t$. Every entry of $A_0$
is less than $3\sqrt p$. Hence
\begin{equation}\label{eq:gap-absolute-remainder}
 0<\mathcal R_l<(2^h-1)3^{h-1}p^{h/2}
 \lambda^{(h-2)l+t}.
\end{equation}
The constant lower-state path is included in this estimate.

\begin{lemma}[Positivity of the leading discrepancy]
\label{lem:positive-discrepancy}
Suppose a positive word of $h\geqslant2$ factors in the model
\eqref{eq:gap-positive-model}, with any integral diagonal gaps, has
normalized trace $u_{\eps m}$, where $m\geqslant1$. Let $X$ denote
either of its constant-state contributions. If
$Y=a\lambda^m$, with $a=c$ for $\eps=1$ and $a=d$ for $\eps=-1$,
then $Y>X$.
\end{lemma}

\begin{proof}
The two constant contributions have product
$C=3^{2h-2}K^h=K(9K)^{h-1}>K$. The normalized word trace therefore
exceeds $X+C/X$. On the other hand the single trace is $Y+K/Y$, and
the minimum ratio gives $Y^2/K\geqslant\lambda^{2m-1}>1$.
If $Y\leqslant X$, monotonicity of $z+K/z$ on $z>\sqrt K$ yields
$$
 V>X+C/X>X+K/X\geqslant Y+K/Y,
$$
a contradiction. The argument applies separately to either constant
path; it does not require that the chosen one dominate the other.
\end{proof}

\subsection{A finite window for the single exponent}

At an equality $V_l=u_{\eps m}$, write
$$
 (a,b)=\begin{cases}(c,d),&\eps=1,\\(d,c),&\eps=-1,
 \end{cases}
 \qquad q=m-hl-t,\qquad H=\lceil h/2\rceil.
$$
The two possible values of $c^h/a$ are
$K^{(h-1)/2}(c/d)^{(h-1)/2}$ and
$K^{(h-1)/2}(c/d)^{(h+1)/2}$. Consequently
\begin{equation}\label{eq:gap-coefficient-window}
 \lambda^{-(h+1)/2}
 <3^{h-1}\frac{c^h}{a}
 \leqslant\left(\frac3{\sqrt5}\right)^{h-1}
 \lambda^{(h+1)/2}.
\end{equation}
Since $m\geqslant1$, the reciprocal term in the single entry gives
\begin{equation}\label{eq:gap-single-leading-ratio}
 1\leqslant\frac{u_{\eps m}}{a\lambda^m}
 \leqslant1+\lambda^{-1}<\frac32.
\end{equation}
Equations \eqref{eq:gap-coefficient-window} and
\eqref{eq:gap-single-leading-ratio}, together with $V_l>X_l$, imply
$$
 \lambda^q>\frac23\lambda^{-(h+1)/2}
 >\lambda^{-(h+2)/2}.
$$
The last inequality follows from $\lambda>5/2>9/4$. Thus the integral
exponent satisfies $q\geqslant-H$.

For $l\geqslant1$, use \eqref{eq:gap-binomial-bound} for the upper
estimate. The inequalities
$1+\lambda^{-1}<7/5$ and $21/(5\sqrt5)<2<\lambda$ give
$\lambda^q<\lambda^{(3h+1)/2}$. For $l=0$, the binomial estimate
has an extra factor $\lambda^h$. Integer rounding, for both parities
of $h$, therefore gives the complete window
\begin{equation}\label{eq:gap-phase-window}
 -\left\lceil\frac h2\right\rceil\leqslant q\leqslant
 \begin{cases}
 \lfloor3h/2\rfloor,&l\geqslant1,\\
 \lfloor5h/2\rfloor,&l=0.
 \end{cases}
\end{equation}

\subsection{The algebraic norm bounds the product index}

Rearrange the equality as
\begin{equation}\label{eq:gap-leading-discrepancy}
 \lambda^{hl+t}\Delta=\mathcal R_l-b\lambda^{-m},
 \qquad \Delta=a\lambda^q-3^{h-1}c^h.
\end{equation}
The two terms defining $\Delta$ have norms, in $\Q(\sqrt D)$,
$K$ and $K(9K)^{h-1}$, respectively. They differ because $9K>1$.
Thus $\Delta\ne0$; this also follows, with its sign, from
Lemma~\ref{lem:positive-discrepancy}. The norm argument additionally
gives a quantitative separation.

The integer $D$ lies strictly between $(3p-1)^2$ and $(3p)^2$, so it
is nonsquare. The number $\lambda$ is an algebraic-integer unit, and
$\sqrt D\,c$, $\sqrt D\,d$ are algebraic integers by the exact ray
formulas. Hence $(\sqrt D)^h\Delta$ is an algebraic integer, whether
or not $D$ is squarefree. Its nonzero integral norm gives
\begin{equation}\label{eq:gap-norm-lower}
 |\Delta\overline\Delta|\geqslant D^{-h}.
\end{equation}
Conjugation exchanges $c,d$ and replaces $\lambda$ by $\lambda^{-1}$.
Using \eqref{eq:gap-phase-window}, $c,d<\sqrt p$, and
$H\leqslant h-1$, we obtain
\begin{align*}
 |\overline\Delta|
 &<3^H p^{H+1/2}+3^{h-1}p^{h/2}\\
 &\leqslant2\cdot3^{h-1}p^{H+1/2}.
\end{align*}
Since $D<9p^2$, equation~\eqref{eq:gap-norm-lower} implies
\begin{equation}\label{eq:gap-discrepancy-lower}
 |\Delta|>
 \frac1{2\cdot3^{3h-1}p^{2h+H+1/2}}.
\end{equation}
Meanwhile \eqref{eq:gap-absolute-remainder}, $b<\sqrt p$ and
$m\geqslant1$ show that the absolute value of the right side of
\eqref{eq:gap-leading-discrepancy} is less than
$2^h3^{h-1}p^{h/2}\lambda^{(h-2)l+t}$. Combining the two estimates
and cancelling positive powers of $\lambda$ gives
\begin{equation}\label{eq:gap-index-inequality}
 \lambda^{2l}
 <2^{h+1}3^{4h-2}p^{5h/2+H+1/2}
 \leqslant2^{h+1}3^{4h-2}p^{3h+1}.
\end{equation}
The last step uses $H\leqslant(h+1)/2$. Now
$$
 2^{h+1}3^{4h-2}=\frac29\,162^h,
 \qquad (5/2)^6=\frac{15625}{64}>162.
$$
If $l\geqslant3h$, the lower bound $\lambda>(5/2)p$ contradicts
\eqref{eq:gap-index-inequality}, since $6h\geqslant3h+1$. Therefore
\begin{equation}\label{eq:gap-product-index-bound}
 l\leqslant3h-1.
\end{equation}
For $l\geqslant1$, equations~\eqref{eq:gap-reconstruction} and
\eqref{eq:gap-phase-window} give
$$
 k\leqslant l+m=(h+1)l+t+q
 \leqslant(h+1)(3h-1)+h-1+\lfloor3h/2\rfloor=b(h).
$$
For $l=0$ they give the smaller bound
$k\leqslant h-1+\lfloor5h/2\rfloor$. This proves
Theorem~\ref{thm:gap-bound}.

The conclusion bounds the gap at fixed $h$, rather than the occurrence
count itself. More importantly for the next sections, it supplies the
complete index domain \eqref{eq:gap-phase-window}--\eqref{eq:gap-product-index-bound}
without imposing a numerical cutoff on the Markov parameter $p$.
The remaining task is to bound that parameter independently of the ray.
